\section{Verification purpose}

This case verifies how firebrand ignition delay alters spread when surface propagation and spotting operate together. The surface-fire-dominated transport simulation condition delays firebrand development until the surface front controls advance; combined surface-fire and firebrand transport permits sufficiently distant deposits to develop ahead of that front.

The paired design exposes incorrect delay application, a missing transition between propagation modes, or an inconsistent surface-rate reference. Success requires surface-controlled behavior in the first simulation condition and quantified earlier arrival, spot ignition, and leading-edge gain in the second.

\section{Mathematical formulation and numerical implementation}

\subsection{Surface-fire and firebrand time scales}

The verification configuration is a flat, uniform fire behavior fuel model (FBFM) 102 grassland with wind
speed $U_w=\SI{6.71}{\metre\per\second}$. The surface head-fire spread rate
$U_s$ is read from the surface-dominated simulation condition spread rate raster produced by the Rothermel
model in the tested ELMFIRE executable. For the reported run,
$U_s=\SI{\Metric{surfacerosreferencemps}}{\metre\per\second}$. In the absence of a successful
firebrand-induced ignition, the exact one-dimensional solution is
\begin{equation}
  t_s(x)=\frac{x}{U_s}, \qquad
  x_{\mathrm{LE},s}(t)=U_s t, \qquad
  \overline{\mathrm{ROS}}_{\mathrm{LE},s}=U_s .
  \label{eq:surface}
\end{equation}

For a firebrand emitted at $x_s$ and transported a downwind distance $X$, the
earliest time at which it can produce a fully developed fire is
\begin{equation}
 t_{\mathrm{spot}}(x_s,X)
 = \frac{x_s}{U_s}+\frac{X}{U_w}
   +t_{\mathrm{ign,small}}+t_{\mathrm{ign,large}} .
 \label{eq:spot-time}
\end{equation}
The surface front reaches the same target location at
$t_s(x_s+X)=(x_s+X)/U_s$. With immediate smoldering-to-flaming transition,
a spot fire can lead the surface front only when
\begin{equation}
 X\left(\frac{1}{U_s}-\frac{1}{U_w}\right)
 > t_{\mathrm{ign,large}} .
 \label{eq:lead-condition}
\end{equation}

\subsection{ELMFIRE implementation exercised}

The current implementation uses the non-superseded spotting path with these
explicit selectors:
\begin{itemize}
\item generation proportional to fire power;
\item an empirical transport-distance distribution; and
\item Eulerian accumulation.
\end{itemize}
For each burning cell, the firebrand-generation calculation evaluates
\begin{equation}
 \dot N_{\mathrm{cell}}
 = \frac{I_f\,\Delta y}{1000}\,GR^* ,
 \label{eq:generation}
\end{equation}
where $I_f$ is fireline intensity in \si{\kilo\watt\per\metre},
$\Delta y$ is the cell width, and
$GR^*=\SI{33.3}{\per\second\per\mega\watt}$. The factor
$I_f\Delta y/1000$ is the heat-release rate of the cell in MW. The Eulerian
transport calculation deposits these particles according to the configured
Sardoy-type distribution ($\mu=2.18$, $\sigma=1.23$), truncated at an
effective maximum distance of approximately \SI{150}{\metre}.

The ignition calculation first converts the deposited number to a number per unit
area. In the surface-dominated simulation condition, the physical ignition model
applies the
wildland mass-load criterion before sampling SFT, with
$\vartheta_{\mathrm{ign,small}}=\SI{42}{\second}$ and
$t_{\mathrm{ign,large}}=\SI{300}{\second}$. For this parameter set, the
documented limiting behavior predicts that deposited mass remains below the
physical threshold, so the SFT probability is zero~\cite{Qin2025}. In the
mixed-mode simulation condition, the simple ignition model
prescribes immediate
SFT and a \SI{150}{\second} development delay. When the development delay
expires, the level-set state, arrival-time field, and ember-ignition diagnostic
are updated by the same firebrand ignition event.

\section{Derivation of expected behavior}

At the maximum transport distance $X_{\max}=\SI{150}{\metre}$, flight takes
\[
 t_f=\frac{150}{6.71}=\SI{22.4}{\second},
 \qquad
 t_s(150)=\frac{150}{0.577}=\SI{260.0}{\second}.
\]
For surface-dominated simulation condition, even the deliberately optimistic assumption of immediate SFT
would give $t_f+t_{\mathrm{ign,large}}=\SI{322.4}{\second}$, later than the
surface arrival. The expected result is therefore Equation~\ref{eq:surface},
with no cells attributed to firebrand ignition.

For mixed-mode simulation condition, the earliest full spot fire at $X_{\max}$ occurs at
$22.4+150=\SI{172.4}{\second}$. Moreover, the left side of
Equation~\ref{eq:lead-condition} is \SI{237.6}{\second}, which exceeds the
\SI{150}{\second} development delay. Long-distance deposits can therefore
mature ahead of the primary front. The expected signature is an acceleration
near \SI{170}{\second}, at least one firebrand-ignited cell, a leading edge
ahead of $U_s t$, and a mean spread rate greater than $U_s$ over the common
analysis interval. These are the qualitative features expected from mixed-mode spread.

% BEGIN EXPLICIT EXPECTED-RESULT DERIVATION
The lead condition can also be solved for the minimum useful spotting
distance:
\[
 X_{\mathrm{crit}}=
 \frac{t_{\mathrm{ign,large}}}{1/U_s-1/U_w},\qquad
 \frac{1}{0.577}-\frac{1}{6.71}=1.584071~\mathrm{s\,m^{-1}}.
\]
For the surface-dominated 300-s delay,
\(X_{\mathrm{crit}}=300/1.584071=189.39~\mathrm{m}\), which exceeds the
configured \(X_{\max}=150~\mathrm{m}\); even the longest deposit cannot lead.
For the mixed 150-s delay,
\(X_{\mathrm{crit}}=94.69~\mathrm{m}<150~\mathrm{m}\), so the tail
\(94.69<X\leq150~\mathrm{m}\) can mature ahead of the surface front.
At \(X=150~\mathrm{m}\),
\[
 t_f=150/6.71=22.3547~\mathrm{s},\quad
 t_s=150/0.577=259.9653~\mathrm{s},\quad
 t_s-t_f=237.6106~\mathrm{s}.
\]
Thus the mixed earliest full-spot scale is
\(22.3547+150=172.3547~\mathrm{s}\), whereas the
surface-dominated optimistic scale is
\(22.3547+300=322.3547~\mathrm{s}>t_s\).
These substitutions yield the explicit reference signatures used by the
metrics: zero surface-dominated ember ignitions, positive mixed-mode ember
ignitions, mixed \(ROS/U_s\geq1.10\), at least 20 m edge and paired-time of arrival (TOA)
advantages, and onset in the 120--220 s interval around 172.35 s.
% END EXPLICIT EXPECTED-RESULT DERIVATION

\section{Verification metrics and acceptance criteria}

The postprocessor excludes the two-cell halo and operates on the centre-row
TOA and spread-rate profiles. It converts the positive median Rothermel
spread rate value from ft/min to m/s and uses it as $U_s$ throughout the
reference calculations. It fits $t=a+x/U$ over
$\SI{50}{\metre}\le x\le\SI{280}{\metre}$ using at least 12 cells.
The following independently informative conditions must all pass:

\begin{enumerate}
\item surface-dominated simulation condition mean TOA relative error against $x/U_s$ is at most 0.10.
\item surface-dominated simulation condition fitted rate of spread (ROS) relative error against $U_s$ is at most 0.10.
\item surface-dominated simulation condition contains no firebrand-ignited physical cells.
\item mixed-mode simulation condition contains at least one firebrand-ignited physical cell.
\item mixed-mode simulation condition fitted ROS is at least $1.10U_s$.
\item At \SI{250}{\second}, its leading edge is at least \SI{20}{\metre}
      ahead of $U_s t$.
\item Its first one-cell leading-edge gain occurs from 120 to
      \SI{220}{\second}, bracketing the derived \SI{172.4}{\second} scale.
\item Its median TOA is at least \SI{20}{\second} earlier than surface-dominated simulation condition over
      the shared fitting interval.
\end{enumerate}

Missing outputs do not pass by default. Both simulation conditions must provide final TOA,
accumulated-firebrand, and ember-ignition GeoTIFFs before any case-level pass
decision is possible.

\subsection{Justification of metric quantification}

The two 10\% surface-baseline limits allow cell-centred TOA and regression
quantization while requiring agreement with the run-derived Rothermel speed
over a \SI{230}{m} fitting interval, more than twenty grid cells.  The zero
spot count in the surface-dominated simulation condition and the positive count in the
mixed simulation condition are exact mechanism checks and therefore have no numerical
tolerance.  For the mixed mode, \(ROS/U_s\ge1.10\) requires a 10\% effect
beyond the baseline fitting band, while the \SI{20}{m} edge and paired-TOA
leads equal two spatial cells and are too large to arise from one-cell front
classification alone.

The 120--\SI{220}{s} onset window brackets the independently derived
\SI{172.4}{s} earliest full-spot scale by approximately \(\pm50\) s.  That
width accommodates stochastic deposition, cell crossing, and the fact that
the metric detects the first full-cell gain rather than a continuous event,
but still rejects acceleration unrelated to the predicted delay regime.  The
eight behavior checks are conjunctive because a faster fitted front without
ember attribution, timing consistency, and a paired TOA lead would not verify
firebrand-delay coupling.

\subsection{Predeclared acceptance criteria}

Table~\ref{tab:acceptance-contract} summarizes the metrics fixed before
inspection of the calculated results. Numerical values and component statuses
are reported separately in Section~7.

\begin{table}[H]
\centering
\footnotesize
\caption{Predeclared verification metrics and acceptance criteria.}
\label{tab:acceptance-contract}
\begin{tabular}{@{}p{0.23\linewidth}p{0.47\linewidth}p{0.21\linewidth}@{}}
\toprule
Metric & Output, region, and aggregation & Acceptance criterion\\
\midrule
Output completeness & Final TOA, accumulated-firebrand, and ember-ignition rasters for both current simulation conditions; halo and nodata excluded. & 2 of 2 simulation conditions \\
Surface-dominated TOA error & Mean relative error against \(x/U_s\) on 50--\SI{280}{m}; at least 12 cells. & \(\le0.10\) \\
Surface-dominated ROS error & Relative error of fitted ROS against run-derived Rothermel \(U_s\). & \(\le0.10\) \\
Surface-dominated spot count & Number of ember-ignited physical cells. & 0 \\
Mixed-mode spot count & Number of ember-ignited physical cells. & \(\ge1\) \\
Mixed-mode ROS ratio & Fitted mixed-mode ROS divided by \(U_s\). & \(\ge1.10\) \\
Mixed-mode edge gain & Leading-edge distance ahead of \(U_st\) at \SI{250}{s}. & \(\ge\SI{20}{m}\) \\
Acceleration onset & First time the leading edge exceeds the surface reference by one cell. & 120--\SI{220}{s} \\
Paired TOA lead & Median surface-dominated minus mixed-mode TOA over common fitting cells. & \(\ge\SI{20}{s}\) \\
Overall decision & Conjunction of completeness and all eight behavior checks. & PASS only if all pass \\
\bottomrule
\end{tabular}
\end{table}

\section{Simulation configuration and predicted outcome}

Figure~\ref{fig:input-configuration} records the prepared spatial inputs for combined surface-fire and firebrand transport.
These panels show the prepared spatial fields, remove the
two-cell numerical buffer, and retain map coordinates in metres.  The left panel shows
the most informative fuel or structure layout available for the case; the right panel
shows the cells selected by the non-positive initial level-set field, making a localized
ignition or firebrand-emission source explicit.

\begin{figure}[H]
\centering
\EvidenceFigure[width=0.98\textwidth,height=0.42\textheight,keepaspectratio]{../figures/input_configuration.pdf}
\caption{Whole-domain prepared input configuration for the representative simulation condition.  The
physical domain is shown without the numerical halo; coordinates, configured wind values,
fuel or structure layout, and initial ignition/source geometry are identified in the panels.}
\label{fig:input-configuration}
\end{figure}

The Courant--Friedrichs--Lewy (CFL) number measures the distance
advected during one timestep relative to the grid spacing and
constrains the time discretization.

\begin{table}[htbp]
\centering
\caption{Parameters common to both verification simulation conditions.}
\label{tab:common-config}
\begin{tabular}{lll}
\toprule
Parameter & Value & Verification role \\
\midrule
Physical domain & $300\times100$ m & one-dimensional centre transect \\
Raster domain & $34\times14$ cells & includes two-cell halo \\
Cell size & \SI{10}{\metre} & controlled baseline resolution \\
Time step / wind CFL & \SI{0.742574}{\second} / 0.4979 & common exact-duration divisor \\
Fuel and terrain & FBFM 102, flat & homogeneous surface spread \\
Wind & \SI{6.71}{\metre\per\second} toward $+x$ & firebrand transport \\
Ignition geometry & full crosswind line at $x=\SI{5}{\metre}$ & 1-D behavior \\
Generation & PER-MW, $GR^*=33.3$ & heat-release-scaled source \\
Transport / accumulation & empirical / Eulerian & reference transport path \\
Surface fire & enabled & general coupled fire spread \\
\bottomrule
\end{tabular}
\end{table}

\begin{table}[htbp]
\centering
\caption{Ignition parameters changed by the paired experiment.}
\label{tab:variant-config}
\begin{tabular}{p{0.27\textwidth}p{0.12\textwidth}p{0.13\textwidth}p{0.13\textwidth}p{0.17\textwidth}}
\toprule
Simulation condition & Ignition model & $t_{\mathrm{ign,small}}$ & $t_{\mathrm{ign,large}}$ & Prediction \\
\midrule
surface-dominated simulation condition & PHYSICAL & \SI{42}{\second} scale & \SI{300}{\second} & surface dominated \\
mixed-mode simulation condition & SIMPLE & \SI{0}{\second} & \SI{150}{\second} & mixed-mode acceleration \\
\bottomrule
\end{tabular}
\end{table}

\subsection{Preprocessing, execution, and postprocessing}

The two simulation conditions use independently prepared, co-registered spatial fields. Both represent flat terrain, fuel model 102, and
constant wind; places a line ignition across every physical crosswind row; and
writes explicit model selectors and ignition parameters into each local
ELMFIRE simulation configuration. The halo is present numerically but is excluded from
all reported quantities. The nominal 0.5-wind-CFL step is 0.745645~s. Because the controlled stops are 300 and 450~s, their greatest common duration is 150~s; preprocessing selects \(N=\lceil150/0.745645\rceil=202\) and the common \(\Delta t=150/202=0.742574\)~s. Both stops are therefore exact timestep multiples and the actual wind CFL remains below 0.5.

Both prescribed conditions are simulated independently with ELMFIRE. Their
final physical fields are compared using the surface-rate reference and the
arrival-time, ignition, and leading-edge measures defined above.

\clearpage
\section{Actual simulation results}

Figure~\ref{fig:domain-result} shows the representative time-of-arrival field from combined surface-fire and firebrand transport.
The displayed field is taken from the simulated results, masks missing values, and excludes
the two-cell numerical buffer.  The domain map complements the spread-rate summaries by showing the combined surface-fire and delayed-firebrand arrival pattern.

\begin{figure}[H]
\centering
\EvidenceFigure[width=0.96\textwidth,height=0.50\textheight,keepaspectratio]{../figures/domain_result.pdf}
\caption{Whole-domain time-of-arrival field from the representative completed ELMFIRE run.  This
domain-scale result supplements, rather than replaces, the higher-level verification
profiles, convergence plots, ensemble statistics, and calculated acceptance metrics below.}
\label{fig:domain-result}
\end{figure}

Figures~\ref{fig:surface-dominated-reference} and
\ref{fig:mixed-mode-reference} first present each ignition regime independently
using the same four observables: accumulated firebrands, arrival time,
leading-edge position, and leading-edge ROS. Colored curves are actual ELMFIRE
outputs; black curves are the Rothermel surface-spread references.

\begin{figure}[H]
\centering
\EvidenceFigure[width=0.96\textwidth,height=0.66\textheight,keepaspectratio]{../figures/surface_dominated_reference.pdf}
\caption{Long-delay surface-dominated simulation condition. The leading edge remains close
to the surface-spread characteristic despite continued firebrand deposition.}
\label{fig:surface-dominated-reference}
\end{figure}

\begin{figure}[H]
\centering
\EvidenceFigure[width=0.96\textwidth,height=0.66\textheight,keepaspectratio]{../figures/mixed_mode_reference.pdf}
\caption{Finite-delay mixed-mode simulation condition. Departure from the surface-only
arrival, trajectory, and ROS references identifies firebrand-induced advance.}
\label{fig:mixed-mode-reference}
\end{figure}

Figure~\ref{fig:delay-effect} overlays the two simulation conditions as a supplemental
controlled comparison, making the timing and magnitude of their divergence
easier to identify.

\begin{figure}[H]
\centering
\EvidenceFigure[width=0.96\textwidth,height=0.66\textheight,keepaspectratio]{../figures/mixed_mode_delay_effect.pdf}
\caption{Paired ELMFIRE comparison of the long-delay surface-dominated and
finite-delay mixed-mode regimes. Accumulation is crosswind averaged; the other
panels use the centre-row arrival-time field.}
\label{fig:delay-effect}
\end{figure}
\FloatBarrier

\section{Calculated metrics and verification decision}

Table~\ref{tab:metrics} reads the calculated entries directly from
calculated verification and validation metrics. An empty simulation result is represented as
N/A with NOT EVALUABLE, never as an empty table cell or an implicit pass.

\begin{table}[H]
\centering
\small
\caption{Acceptance metrics calculated from the paired ELMFIRE outputs.}
\label{tab:metrics}
\begin{tabular}{llll}
\toprule
Metric & Acceptance limit & Calculated & Status \\
\midrule
Output completeness & 2 simulation conditions & \Metric{completedvariants}/2 & \Metric{status} \\
surface-dominated simulation condition mean TOA error & $\le 0.10$ & \Metric{surfacedominatedtoameanrelativeerror} & \Metric{surfacedominatedtoastatus} \\
surface-dominated simulation condition mean ROS (\si{\metre\per\second}) & $\Metric{surfacerosreferencemps}\mathbin{\pm}10\%$ & \Metric{surfacedominatedmeanrosmps} & \Metric{surfacedominatedrosstatus} \\
surface-dominated simulation condition ember-ignited cells & $\le 0$ & \Metric{surfacedominatedemberignitioncells} & \Metric{surfacedominatedignitionstatus} \\
mixed-mode simulation condition ROS / $U_s$ & $\ge 1.10$ & \Metric{mixedmoderosratiotosurface} & \Metric{mixedmoderosstatus} \\
mixed-mode simulation condition ember-ignited cells & $\ge 1$ & \Metric{mixedmodeemberignitioncells} & \Metric{mixedmodeignitionstatus} \\
mixed-mode simulation condition edge gain at 250 s (m) & $\ge 20$ & \Metric{mixedmodeleadingedgegainm} & \Metric{mixedmodeedgegainstatus} \\
mixed-mode simulation condition acceleration onset (s) & $120$--$220$ & \Metric{mixedmodeaccelerationonsets} & \Metric{mixedmodeonsetstatus} \\
Median mixed-mode simulation condition TOA lead (s) & $\ge 20$ & \Metric{pairedmediantoaleads} & \Metric{pairedtoastatus} \\
\midrule
Overall decision & all components pass & \Metric{verificationpassed} & \Metric{status} \\
\bottomrule
\end{tabular}
\end{table}

\section{Technical interpretation}

The current paired execution passes every component. The surface-dominated mean ROS is \Metric{surfacedominatedmeanrosmps}~m/s with no ember-ignited cells. The mixed-mode mean ROS is \Metric{mixedmodemeanrosmps}~m/s (\Metric{mixedmoderosratiotosurface} times the surface reference), with \Metric{mixedmodeemberignitioncells} ember-ignited cells, a \Metric{mixedmodeleadingedgegainm}-m leading-edge gain, acceleration onset at \Metric{mixedmodeaccelerationonsets}~s, and a median TOA lead of \Metric{pairedmediantoaleads}~s.

A pass demonstrates two complementary properties: delayed firebrand states do
not spuriously outrun the surface solver in the surface-dominated limiting
case, and reducing the delay activates a resolvable acceleration in the
mixed-mode simulation condition.
If surface-dominated simulation condition fails alone, inspect the physical mass threshold, development-state
timing, and arrival-time update. If mixed-mode simulation condition fails alone, inspect the simple
ignition probability, persistence of deposited ember states across level-set
steps, delayed transition to a full ignition, and the coupling that writes
that ignition to the level set.

This is numerical verification of the controlled formulation, not field validation.
The ELMFIRE experiment is two-dimensional but deliberately uses a
uniform crosswind line so the centre transect approximates the corresponding
one-dimensional calculation. The mixed-mode thresholds diagnose the expected
mixed regime rather than an unavailable closed-form mixed-mode solution.
Stochastic ignition means the exact onset may vary, which is why its
acceptance window brackets rather than equals \SI{172.4}{\second}.

The current analysis note is: \Metric{notes}

\section{Scope and limitations}
This paired test isolates the interaction between surface propagation and firebrand ignition delay in a uniform domain. It verifies implementation behavior and timing attribution; it does not validate either spread mechanism against field observations.

\begin{thebibliography}{9}
\bibitem{Qin2025}
Y. Qin, \emph{A Physics-Based Eulerian Framework for Modeling Firebrand
Showering in Regional-Scale Wildland and Wildland--Urban Interface Fire
Simulations}, Ph.D. dissertation, University of Maryland, College Park, 2025.
\end{thebibliography}
